%=====================================================================
% KUET CSE 3200 – System Development Project Report
% Project : Housely – A Full-Stack House Rental & Sale Platform
% Compiler : pdflatex  (Overleaf default – do NOT switch to XeLaTeX)
% HOW TO COMPILE:
%   1. Upload kuet_logo.png to your Overleaf project root
%   2. Upload any figure PNGs (arch.png, er.png, etc.) to root
%   3. Click Recompile (pdflatex, twice for TOC/LOF/LOT)
%=====================================================================
\documentclass[12pt,a4paper]{report}

%--------------------------------------------------------------------
% PACKAGES
%--------------------------------------------------------------------
\usepackage[T1]{fontenc}
\usepackage[utf8]{inputenc}
\usepackage{mathptmx}           % Times New Roman – pdflatex compatible
\usepackage[top=1.2in, bottom=1.0in, left=1.2in, right=1.0in,
            headheight=14pt, headsep=0.25in, footskip=0.4in]{geometry}
\usepackage{setspace}
\usepackage{titlesec}
\usepackage{tocloft}
\usepackage{graphicx}
\usepackage{booktabs}
\usepackage{longtable}
\usepackage{array}
\usepackage{caption}
\usepackage{fancyhdr}
\usepackage{hyperref}
\usepackage{enumitem}
\usepackage{xcolor}
\usepackage{tabularx}
\usepackage{multirow}
\usepackage{amsmath}
\usepackage{microtype}
\usepackage{url}
\usepackage{listings}           % code listings
\usepackage{tikz}               % state-machine diagrams
\usetikzlibrary{automata, positioning, arrows.meta, shapes.geometric}
\usepackage{pgfplots}           % bar chart for test results
\pgfplotsset{compat=1.18}
\usepackage{pgfgantt}           % Gantt chart
\usepackage{float}              % [H] placement

%--------------------------------------------------------------------
% ROMAN CHAPTER NUMBERS
%--------------------------------------------------------------------
\renewcommand{\thechapter}{\Roman{chapter}}

%--------------------------------------------------------------------
% LINE SPACING
%--------------------------------------------------------------------
\setstretch{1.5}

%--------------------------------------------------------------------
% CAPTION STYLE  (11pt, centered; table caption above)
%--------------------------------------------------------------------
\captionsetup{font=small, labelfont=bf, justification=centering,
              skip=6pt, position=below}
\captionsetup[table]{position=above}

%--------------------------------------------------------------------
% CHAPTER HEADING – 16pt, Bold, Centered
%--------------------------------------------------------------------
\titleformat{\chapter}[block]
  {\fontsize{16}{24}\bfseries\centering}
  {CHAPTER \thechapter}{1em}{}
\titlespacing*{\chapter}{0pt}{0pt}{24pt}

%--------------------------------------------------------------------
% SECTION – 14pt, Bold, Left
%--------------------------------------------------------------------
\titleformat{\section}
  {\fontsize{14}{21}\bfseries\raggedright}
  {\thesection}{1em}{}
\titlespacing*{\section}{0pt}{12pt}{6pt}

%--------------------------------------------------------------------
% SUBSECTION – 12pt, Bold, Left
%--------------------------------------------------------------------
\titleformat{\subsection}
  {\fontsize{12}{18}\bfseries\raggedright}
  {\thesubsection}{1em}{}
\titlespacing*{\subsection}{0pt}{12pt}{6pt}

%--------------------------------------------------------------------
% HEADER / FOOTER
%--------------------------------------------------------------------
\pagestyle{fancy}
\fancyhf{}
\fancyhead[R]{\small CSE 3200: System Development Project}
\fancyfoot[C]{\thepage}
\renewcommand{\headrulewidth}{0.4pt}
\fancypagestyle{plain}{%
  \fancyhf{}
  \fancyfoot[C]{\thepage}
  \renewcommand{\headrulewidth}{0pt}}

%--------------------------------------------------------------------
% TOC TITLES
%--------------------------------------------------------------------
\renewcommand{\contentsname}{Contents}
\renewcommand{\listtablename}{List of Tables}
\renewcommand{\listfigurename}{List of Figures}

%--------------------------------------------------------------------
% PARAGRAPH SPACING
%--------------------------------------------------------------------
\setlength{\parskip}{6pt}
\setlength{\parindent}{0pt}

%--------------------------------------------------------------------
% CODE LISTING STYLE
%--------------------------------------------------------------------
\definecolor{codegray}{rgb}{0.5,0.5,0.5}
\definecolor{codebg}{rgb}{0.97,0.97,0.97}
\lstdefinestyle{jsStyle}{
  backgroundcolor=\color{codebg},
  commentstyle=\color{codegray}\itshape,
  keywordstyle=\bfseries,
  numberstyle=\tiny\color{codegray},
  stringstyle=\color{black},
  basicstyle=\ttfamily\small,
  breakatwhitespace=false,
  breaklines=true,
  keepspaces=true,
  numbers=left,
  numbersep=6pt,
  showstringspaces=false,
  tabsize=2,
  frame=single,
  framesep=4pt,
  xleftmargin=18pt
}
\lstset{style=jsStyle}

%--------------------------------------------------------------------
% HYPERREF
%--------------------------------------------------------------------
\hypersetup{
  colorlinks=true, linkcolor=black, citecolor=black, urlcolor=blue,
  pdftitle={Housely -- A Full-Stack House Rental and Sale Platform},
  pdfauthor={MD Jahid Hasan Jim \& Shakhoyat Rahman Shujon}}

%====================================================================
\begin{document}
%====================================================================


%--------------------------------------------------------------------
% COVER PAGE
%--------------------------------------------------------------------
\begin{titlepage}
\thispagestyle{empty}
\begin{center}

{\fontsize{14}{21}\normalfont CSE 3200: System Development Project}

\vspace{6pt}

{\fontsize{18}{27}\bfseries\scshape
Housely: A Full-Stack House Rental\\[4pt]
and Sale Platform}

\vspace{24pt}

{\fontsize{14}{21}\normalfont By}

\vspace{36pt}

{\fontsize{14}{21}\bfseries MD Jahid Hasan Jim}\\[4pt]
{\fontsize{14}{21}\normalfont Roll: 2107054}

\vspace{8pt}
{\fontsize{14}{21}\bfseries \&}
\vspace{8pt}

{\fontsize{14}{21}\bfseries Shakhoyat Rahman Shujon}\\[4pt]
{\fontsize{14}{21}\normalfont Roll: 2107104}

\vspace{48pt}

\includegraphics[width=1.14in, height=1.5in]{kuet.png}

\vspace{24pt}

\vfill

{\fontsize{12}{18}\bfseries
Department of Computer Science and Engineering\\
Khulna University of Engineering \& Technology\\
Khulna 9203, Bangladesh\\
April, 2026}

\end{center}
\end{titlepage}

%--------------------------------------------------------------------
% TITLE PAGE  (inner page – roman page i)
%--------------------------------------------------------------------
\pagenumbering{roman}
\setcounter{page}{1}
\thispagestyle{plain}

\begin{center}
{\fontsize{18}{27}\bfseries
Housely: A Full-Stack House Rental\\[4pt]
and Sale Platform}

\vspace{36pt}

{\fontsize{12}{18}\normalfont By}

\vspace{24pt}

{\fontsize{12}{18}\bfseries MD Jahid Hasan Jim}\\
{\fontsize{12}{18}\normalfont Roll: 2107054}\\[6pt]
{\fontsize{12}{18}\bfseries \&}\\[6pt]
{\fontsize{12}{18}\bfseries Shakhoyat Rahman Shujon}\\
{\fontsize{12}{18}\normalfont Roll: 2107104}

\vspace{48pt}
\end{center}

\noindent{\fontsize{12}{18}\bfseries Supervisor:}

\vspace{4pt}
\hspace{0.8in}{\fontsize{12}{18}\bfseries Dr. Md. Milon Islam}\\
\hspace{0.8in}{\fontsize{12}{18}\normalfont Assistant Professor}\\
\hspace{0.8in}{\fontsize{12}{18}\normalfont Department of Computer Science and Engineering}\\
\hspace{0.8in}{\fontsize{12}{18}\normalfont Khulna University of Engineering \& Technology}
\hfill\underline{\hspace{1.5in}}

\hfill Signature

\vspace{36pt}
\begin{center}
{\fontsize{12}{18}\normalfont
Department of Computer Science and Engineering\\
Khulna University of Engineering \& Technology\\
Khulna 9203, Bangladesh}

\vspace{6pt}
{\fontsize{16}{24}\bfseries April, 2026}
\end{center}
\clearpage

%--------------------------------------------------------------------
% ACKNOWLEDGMENT
%--------------------------------------------------------------------
\chapter*{Acknowledgment}
\addcontentsline{toc}{chapter}{Acknowledgment}
\thispagestyle{plain}

All praise goes to the Almighty Allah, whose blessings and mercy enabled us to bring this
project to a successful completion. We extend our deepest gratitude to our supervisor,
\textbf{Dr. Md. Milon Islam}, Assistant Professor, Department of Computer Science and
Engineering, Khulna University of Engineering \& Technology, for his invaluable guidance,
continuous encouragement, and thoughtful feedback at every stage of this work.

We are equally thankful to the faculty members of the Department of Computer Science
and Engineering at KUET for their ongoing support and for providing us with the
academic foundation necessary to undertake a project of this scale. We also acknowledge
the open-source communities behind React Native, Next.js, Express.js, Prisma, and
Clerk, whose well-maintained documentation and tooling made much of this work
possible. Lastly, we owe sincere thanks to our families and friends for their patience and
moral support throughout the duration of this project.

\vspace{24pt}
\hfill{\fontsize{12}{18}\bfseries Authors}

\clearpage

%--------------------------------------------------------------------
% ABSTRACT
%--------------------------------------------------------------------
\chapter*{Abstract}
\addcontentsline{toc}{chapter}{Abstract}
\thispagestyle{plain}

The property rental and sale market in Bangladesh remains largely fragmented, depending
on informal telephone channels and unverified broker networks that offer little transparency
or accountability. This project addresses that gap with \textbf{Housely}: a multi-platform,
integrated house rental and sale coordination system comprising a React Native mobile
application, a Node.js/Express.js REST API backend, and a Next.js administration
dashboard.

The system supports three distinct actor roles---USER, AGENT, and ADMIN---each with a
role-scoped interface secured by Clerk, a third-party authentication and identity
management service. Core operational workflows follow a structured booking lifecycle that
transitions through four states from PENDING to COMPLETED, enforced by a service-layer
finite-state machine (FSM). Real-time in-app messaging between users and agents is
implemented through Socket.IO WebSocket connections authenticated via Clerk session
tokens at the handshake level. Media assets including property images and video tours are
managed through Cloudinary, mobile push notifications are dispatched through Firebase
Cloud Messaging (FCM), and transactional email is delivered through Nodemailer.

The mobile application covers the complete user and agent journey: onboarding, property
discovery with advanced spatial and attribute filtering, booking management, in-app
messaging, and real-time notifications. The admin dashboard provides platform-level
analytics, listing moderation, user management, and booking oversight using interactive
Recharts visualisations.

The platform was built over sixteen weeks following an iterative engineering process
that spanned requirements analysis, database schema design, backend module development,
frontend implementation, and integration testing using Jest and Supertest. The resulting
prototype demonstrates that a modular, type-safe, multi-platform architecture can serve as
a scalable and maintainable framework for digital property coordination, reducing
administrative overhead and improving transparency among property seekers, agents, and
platform administrators.

\clearpage

%--------------------------------------------------------------------
% CONTENTS
%--------------------------------------------------------------------
\tableofcontents
\clearpage

%--------------------------------------------------------------------
% LIST OF TABLES
%--------------------------------------------------------------------
\listoftables
\clearpage

%--------------------------------------------------------------------
% LIST OF FIGURES
%--------------------------------------------------------------------
\listoffigures
\clearpage

%====================================================================
% BODY — Arabic numerals
%====================================================================
\pagenumbering{arabic}
\setcounter{page}{1}

%====================================================================
\chapter{Introduction}
%====================================================================

\section{Introduction}

Rapid urbanisation across Bangladesh has generated considerable demand for residential
property rental and sale, yet the market continues to operate in a largely informal manner.
Prospective tenants and buyers still rely on fragmented classified-advertisement sites,
unverified broker networks, and telephone-based coordination, all of which introduce
unnecessary friction, information asymmetry, and a heightened risk of fraud
\cite{hossain2019}. No dominant digital platform currently exists that simultaneously
aggregates verified listings with multimedia content, enables direct user-to-agent booking
negotiation, supports real-time messaging, and handles post-booking reviews with
administrative moderation.

Housely is designed to fill precisely this gap. Built as a three-component ecosystem
comprising a React Native mobile application, an Express.js REST API, and a Next.js
administration dashboard, the platform digitalises the full property rental and sale journey
for all three actors: property seekers (\textbf{USER}), property owners and agents
(\textbf{AGENT}), and platform administrators (\textbf{ADMIN}).

\section{Background and Problem Statement}

Current approaches to property search in Bangladesh share several structural weaknesses.
Listings are scattered across multiple sites with no unified booking or messaging layer.
Agents have no single interface from which they can receive booking requests, negotiate
terms, and track appointment histories. Users cannot review agents or properties after a
completed transaction, which removes accountability from the process. Platform
administrators lack real-time visibility into listing quality or booking volumes. Housely
addresses all four of these deficiencies within a single cohesive system, with role-scoped
interfaces enforced at both the API and UI layers \cite{bproperty2023, kleppmann2017}.

\section{Objectives}

The principal objectives of this project are as follows.

\begin{enumerate}[noitemsep]
  \item To design and implement a multi-role platform (USER, AGENT, ADMIN) that
        coordinates property discovery, booking, messaging, and review workflows.
  \item To build a secure, scalable REST API backend using Clerk for identity management,
        Redis for rate limiting, and Socket.IO for real-time WebSocket communication.
  \item To deliver usable multi-platform interfaces---a React Native mobile application and
        a Next.js administration dashboard---for workflow execution and business analytics.
  \item To integrate modern cloud services (Cloudinary, Neon PostgreSQL, Firebase FCM,
        and Clerk) into a cohesive and deployable architecture.
  \item To validate the implementation through automated integration testing using Jest
        and Supertest.
\end{enumerate}

\section{Scope}

The project is divided into three subsystems. The \textbf{backend} is built on Node.js~20
(ESM) with Express~4, Prisma~6 ORM, PostgreSQL (Neon serverless), Redis~7 (rate
limiting only), Socket.IO~4, Zod~3 for validation, Multer for file uploads, and Cloudinary
for media storage. Authentication is delegated entirely to Clerk via the
\texttt{@clerk/express} SDK. The backend is deployed as serverless functions on Vercel.
The \textbf{mobile application} is built with React Native~0.81 and Expo~54, using
Expo Router for file-based navigation, NativeWind~4 for Tailwind-based styling,
Zustand~5 for global state, TanStack Query~5 for server-state management, and
\texttt{@clerk/clerk-expo} with \texttt{expo-secure-store} for authentication. The
\textbf{administration dashboard} is built with Next.js~14 (App Router),
\texttt{@clerk/nextjs} for authentication, TanStack React Query~5, Zustand~5, Recharts,
Shadcn/ui, React Hook Form, Zod~4, and Tailwind CSS. Version control was managed
throughout using Git and GitHub.

\section{Unfamiliarity of the Problem and Solution Space}

Established international platforms such as Zillow and Rightmove address property listing
and discovery but are region-specific and do not accommodate Bangladeshi administrative
geography (Division, District, Thana), Taka-denominated pricing, or the three-actor role
model required by local market dynamics. The closest domestic alternative, Bproperty,
offers a web portal but lacks real-time messaging, agent KYC verification, post-booking
reviews, and a unified admin analytics dashboard. Housely's combination of a structured
booking lifecycle FSM, Clerk-managed federated authentication, Socket.IO real-time
messaging with Clerk-authenticated WebSocket handshakes, Firebase FCM push
notifications, and layered React Native navigation represents an integration approach not
found in the reviewed open-source or academic literature \cite{richardson2007, brown2015}.

\section{Project Planning and Work Distribution}

The project was carried out over sixteen weeks by two students. Table~\ref{tab:raci}
presents the RACI matrix for task ownership, and Figure~\ref{fig:gantt} shows the overall
project timeline.

\begin{table}[H]
\caption{RACI Matrix for Project Work Distribution}
\label{tab:raci}
\centering
\small
\setstretch{1.2}
\begin{tabularx}{\textwidth}{|X|c|c|c|c|}
\hline
\textbf{Task} & \textbf{Student~1} & \textbf{Student~2} & \textbf{Supervisor} & \textbf{Reviewer} \\
\hline
Requirements Analysis    & R/A & C   & C & I \\
\hline
DB Schema \& Migrations  & C   & R/A & C & I \\
\hline
Auth \& Clerk Integration & R/A & C   & I & I \\
\hline
House \& Filter Modules  & R/A & C   & I & I \\
\hline
Booking \& Review Modules& C   & R/A & I & I \\
\hline
Message \& Notification  & C   & R/A & I & I \\
\hline
Admin Module             & C   & R/A & I & C \\
\hline
React Native Mobile App  & R/A & C   & I & C \\
\hline
Next.js Admin Dashboard  & C   & R/A & I & C \\
\hline
Testing                  & R   & R/A & I & C \\
\hline
Report Writing           & R/A & R/A & C & C \\
\hline
\end{tabularx}
\end{table}
{\small R\,=\,Responsible, A\,=\,Accountable, C\,=\,Consulted, I\,=\,Informed}
\setstretch{1.5}

\begin{figure}[H]
\centering
\setstretch{1.0}
\begin{ganttchart}[
  hgrid, vgrid,
  x unit=0.52cm, y unit title=0.45cm, y unit chart=0.45cm,
  title height=1, bar height=0.6,
  bar label font=\small, title label font=\small,
  bar/.append style={fill=gray!40}
]{1}{16}
  \gantttitle{Project Timeline (Weeks)}{16} \\
  \gantttitlelist{1,...,16}{1} \\
  \ganttbar{Requirements \& Schema}{1}{2}  \\
  \ganttbar{Auth \& Clerk Integration}{2}{4}\\
  \ganttbar{House \& Filter Modules}{3}{6} \\
  \ganttbar{Booking, Review, Message}{5}{8}\\
  \ganttbar{Notifications \& Admin}{7}{10} \\
  \ganttbar{React Native Mobile App}{5}{12}\\
  \ganttbar{Next.js Admin Dashboard}{7}{12}\\
  \ganttbar{Integration \& Testing}{11}{14}\\
  \ganttbar{Report Writing}{13}{15}        \\
  \ganttbar{Final Presentation}{16}{16}
\end{ganttchart}
\setstretch{1.5}
\caption{Gantt Chart of Project Timeline}
\label{fig:gantt}
\end{figure}

\section{Applications of the Work}

Housely can be deployed as a commercial property listing platform across Bangladesh,
serving major urban centres such as Dhaka, Chittagong, and Khulna, as well as secondary
cities. Beyond its direct commercial application, the architecture functions as a reference
implementation for multi-role marketplace SaaS platforms that require federated identity
management via Clerk, Socket.IO real-time communication, and Cloudinary media
pipelines. The admin analytics layer also provides a reusable template for business
intelligence dashboards in similar marketplace domains.

\section{Organization of the Report}

Chapter~\hyperref[chap:related]{II} reviews related work on digital property platforms,
real-time messaging systems, and cross-platform mobile frameworks.
Chapter~\hyperref[chap:design]{III} describes the system architecture, database schema,
and key design decisions across all three codebases.
Chapter~\hyperref[chap:impl]{IV} presents implementation details, API design, frontend
structure, and test results.
Chapter~\hyperref[chap:societal]{V} addresses societal, ethical, legal, and environmental
considerations.
Chapter~\hyperref[chap:complex]{VI} analyses the complex engineering problems and
activities encountered during the project.
Chapter~\hyperref[chap:conc]{VII} summarises the work, acknowledges limitations, and
outlines directions for future development.

%====================================================================
\chapter{Related Works}
\label{chap:related}
%====================================================================

\section{Introduction}

Property listing and rental platforms span a wide spectrum, ranging from general-purpose
booking APIs to region-specific real-estate marketplaces and purpose-built coordination
systems. This chapter situates Housely within that landscape and identifies the gaps it
is designed to address.

\section{Related Works}

\subsection{International Property Listing Platforms}

Zillow (USA) is the world's largest property marketplace, offering listing search, automated
valuation estimates, and agent directories. Its architecture and user experience are
optimised for the North American market, with USD pricing and US administrative
geography, making it unsuitable for direct adoption in Bangladesh \cite{zillow2023}.
Rightmove (UK) aggregates listings from registered estate agents but does not provide
direct in-app buyer-to-agent messaging, booking management, or any post-listing review
mechanism.

\subsection{Domestic Solutions}

Bproperty is the closest existing commercial solution for the Bangladeshi property market,
offering a web portal and a basic mobile application. However, it does not support
real-time messaging, agent KYC verification, post-booking reviews, a booking lifecycle
state machine, or a unified admin analytics dashboard. Listing management is handled
internally by the Bproperty team rather than by agents directly, which limits platform
scalability.

\subsection{Real-Time Communication Systems}

The use of Socket.IO for real-time messaging in web and mobile systems is well established
in the literature \cite{brown2015}. Marketplace platforms that employ Socket.IO with
token-authenticated handshakes have demonstrated sub-100~ms message delivery latency
in LAN environments, consistent with the performance expectations of Housely's in-app
chat feature.

\subsection{Cross-Platform Mobile Development}

React Native enables code sharing between iOS and Android, reducing development effort
compared to fully native implementations \cite{eisenman2015}. Expo's managed workflow
minimises native build configuration overhead, allowing development teams to focus on
feature delivery. NativeWind 4 brings Tailwind CSS utility classes to React Native,
significantly accelerating the UI implementation process.

\subsection{ORM and Serverless Database Systems}

Prisma ORM has seen wide adoption in the Node.js ecosystem because of its type-safe
query builder and schema-first migration system \cite{prisma2023}. Neon's serverless
PostgreSQL provides autoscaling connection pooling and database branching, which lowers
infrastructure costs for development-stage deployments considerably.

\subsection{Federated Identity Management}

Third-party identity providers such as Clerk consolidate authentication concerns---including
session management, token refresh, social sign-in, and multi-factor authentication---into
a single managed service. Delegating authentication to a purpose-built provider eliminates
entire classes of security vulnerabilities (credential stuffing, token leakage, insecure
storage) and reduces the authentication code surface maintained by the application team.

\section{Discussion}

Table~\ref{tab:related} provides a structured comparison of Housely against the most
relevant existing systems. The comparison confirms that none of the reviewed solutions
covers the full combination of Bangladeshi market specificity, real-time messaging, agent
KYC verification, booking lifecycle management, post-booking reviews, and a unified admin
analytics dashboard. This validates the originality and practical relevance of the Housely
problem formulation.

\begin{table}[H]
\caption{Comparison of Housely with Related Property Platforms}
\label{tab:related}
\centering
\small
\setstretch{1.2}
\begin{tabularx}{\textwidth}{|l|c|c|c|c|c|}
\hline
\textbf{Feature} & \textbf{Housely} & \textbf{Bproperty} & \textbf{Zillow} & \textbf{Rightmove} & \textbf{Lamudi BD} \\
\hline
Bangladesh-specific   & \checkmark & \checkmark & \texttimes & \texttimes & \checkmark \\
\hline
Real-time messaging   & \checkmark & \texttimes & \texttimes & \texttimes & \texttimes \\
\hline
Agent KYC             & \checkmark & Partial    & N/A        & N/A        & \texttimes \\
\hline
Booking lifecycle FSM & \checkmark & \texttimes & \texttimes & \texttimes & \texttimes \\
\hline
Post-booking reviews  & \checkmark & \texttimes & \checkmark & \texttimes & \texttimes \\
\hline
Admin analytics       & \checkmark & Internal   & \checkmark & \checkmark & \texttimes \\
\hline
React Native mobile   & \checkmark & \checkmark & \checkmark & \checkmark & \checkmark \\
\hline
Open-source arch.     & \checkmark & \texttimes & \texttimes & \texttimes & \texttimes \\
\hline
\end{tabularx}
\setstretch{1.5}
\end{table}

%====================================================================
\chapter{System Analysis and Design}
\label{chap:design}
%====================================================================

\section{Introduction}

This chapter describes the methodology, system architecture, module structure, booking
lifecycle state machine, and database schema design decisions across all three codebases
of the Housely platform.

\section{System Analysis and Design Approach}

Development was carried out in five iterative phases.

\begin{enumerate}[noitemsep]
  \item \textbf{Phase 1 -- Requirements and Schema:} Stakeholder analysis, role
        identification, booking lifecycle FSM design, and Prisma schema definition with an
        incremental migration strategy.
  \item \textbf{Phase 2 -- Backend Core:} Express module scaffolding, Clerk
        authentication integration, RBAC middleware, and house listing CRUD operations.
  \item \textbf{Phase 3 -- Backend Features:} Booking, review, message, notification,
        location, and admin analytics modules.
  \item \textbf{Phase 4 -- Frontend Development:} React Native mobile application and
        Next.js admin dashboard, built in parallel against the stabilised API.
  \item \textbf{Phase 5 -- Testing and Documentation:} Jest/Supertest integration tests,
        manual QA across all three platforms, and report writing.
\end{enumerate}

\section{Overall Framework and Architecture}

Figure~\ref{fig:arch} illustrates the overall Housely system architecture. The system
follows a three-tier design: the \textbf{Presentation Tier} contains the React Native
mobile app and the Next.js admin dashboard; the \textbf{Application Tier} is the
Express.js API server responsible for routing, Clerk-based authentication verification,
business logic, and Socket.IO connections; the \textbf{Data Tier} comprises Neon
PostgreSQL (accessed via Prisma~6), Redis~7 (rate limiting only), and Cloudinary (media
storage). The external services layer includes Clerk (identity management), Firebase FCM
(push notifications), and Nodemailer/Gmail (transactional email).

\begin{figure}[H]
\centering
\fbox{\parbox{0.88\textwidth}{\centering\vspace{1.5cm}
\textbf{[Replace with actual System Architecture Diagram]}\\[6pt]
React Native App (\texttt{@clerk/clerk-expo}) \hspace{0.3cm}
Next.js Admin (\texttt{@clerk/nextjs})\\
$\updownarrow$ REST + Socket.IO \hspace{0.5cm} $\updownarrow$ REST\\[4pt]
Express.js API (Node.js 20, Vercel Serverless)\\
$\updownarrow$ Prisma ORM \hspace{0.3cm} $\updownarrow$ ioredis \hspace{0.3cm}
$\updownarrow$ Cloudinary SDK\\[4pt]
Neon PostgreSQL \hspace{0.3cm} Redis 7 (rate limit) \hspace{0.3cm} Cloudinary CDN\\[4pt]
External: Clerk (auth) $|$ Firebase FCM (push) $|$ Nodemailer/Gmail (email)
\vspace{1.5cm}}}
\caption{Overall System Architecture of Housely}
\label{fig:arch}
\end{figure}

\subsection{Data Flow}

Figure~\ref{fig:dfd} presents the Level-1 data flow diagram of Housely, showing the
interactions among the three external actors, the major API processes, persistent data
stores, and external cloud services.

\begin{figure}[H]
\centering
% Replace with \includegraphics{dfd.png} once the diagram image is ready
\fbox{\parbox{0.88\textwidth}{\centering\vspace{1.5cm}
\textbf{[Replace with actual Level-1 Data Flow Diagram]}
\vspace{1.5cm}}}
\caption{Level-1 Data Flow Diagram of Housely}
\label{fig:dfd}
\end{figure}

\subsection{Backend Module Structure}

The Express.js application is organised into ten feature modules, each encapsulating its
own router, controller, service layer, and validation schema. Table~\ref{tab:modules}
summarises these modules along with their primary responsibilities.

\begin{table}[H]
\caption{Express.js Backend Module Structure}
\label{tab:modules}
\centering
\small
\setstretch{1.2}
\begin{tabularx}{\textwidth}{|l|X|}
\hline
\textbf{Module} & \textbf{Responsibility} \\
\hline
\texttt{auth} & Provides a single \texttt{POST /api/auth/sync} endpoint that upserts a Clerk-authenticated user into the local database; the \texttt{protect} and \texttt{requireRole} middlewares delegate identity verification to \texttt{@clerk/express}. \\
\hline
\texttt{user} & User profile retrieval and update, avatar upload to Cloudinary, notification preference management, device token registration for FCM, saved locations, and payment history. \\
\hline
\texttt{house} & Listing CRUD with image and video upload to Cloudinary, favourites, property-view tracking, and discovery endpoints (popular, nearby, recommended, top locations). \\
\hline
\texttt{filter} & Advanced property search accepting up to 12 simultaneous filter parameters, dynamic Prisma \texttt{where} clause construction, and cursor-based pagination. \\
\hline
\texttt{booking} & Booking lifecycle management, four-state FSM transition enforcement, double-booking prevention, and cancellation handling. \\
\hline
\texttt{review} & Post-booking review submission (rating, comment, and media), update, deletion, aggregate rating computation, and listing review display. \\
\hline
\texttt{message} & Conversation creation (one per user-agent pair), message persistence with type support (text/image/audio/video), unread-count tracking, and Socket.IO real-time delivery. \\
\hline
\texttt{notification} & Notification creation, listing, and mark-read operations; Firebase FCM dispatch to registered device tokens; Socket.IO real-time push to authenticated user rooms. \\
\hline
\texttt{location} & Reverse geocoding, saved user location CRUD, and location-based property proximity queries. \\
\hline
\texttt{admin} & Platform analytics (aggregated Prisma queries), listing moderation, user management, and booking oversight. \\
\hline
\end{tabularx}
\setstretch{1.5}
\end{table}

\subsection{Mobile Application Architecture}

The React Native application uses Expo Router for file-based navigation across six
navigator groups: \texttt{(auth)}, \texttt{(onboarding)}, \texttt{(kyc)},
\texttt{(location)}, \texttt{(tabs)} for the user home portal, and \texttt{(owner)} for
the agent portal. Zustand~5 manages global authentication and location state, while
TanStack Query~5 handles server state caching and background refetching. The Axios
instance in \texttt{services/api.js} attaches a fresh Clerk session token on every
outgoing request through a request interceptor that calls the Clerk \texttt{getToken()}
method. Token storage is handled by \texttt{expo-secure-store}, which uses the
platform's secure hardware-backed keychain.

\subsection{Admin Dashboard Architecture}

The Next.js~14 admin dashboard uses the App Router, with server components for initial
data loading and client components for interactive charts. \texttt{@clerk/nextjs} and
\texttt{clerkMiddleware} protect all dashboard routes; unauthenticated requests are
redirected to \texttt{/login}. TanStack React Query~5 handles client-side cache
invalidation, and Recharts renders analytics visualisations for bookings, revenue, user
growth, and listing trends.

\section{Problem Design and Analysis}

\subsection{Role Identification}

Three actor types were identified through property marketplace workflow analysis.

\begin{itemize}[noitemsep]
  \item \textbf{USER} -- Registers an account via Clerk, browses verified property
        listings, applies advanced filters, sends booking requests, exchanges messages with
        agents, submits post-booking reviews, and receives real-time push and in-app
        notifications.
  \item \textbf{AGENT} -- Creates and manages property listings with multimedia uploads,
        confirms or cancels incoming booking requests, communicates with users via in-app
        messaging, and monitors listing statistics and booking queues.
  \item \textbf{ADMIN} -- Exercises system-wide oversight: moderates listings, manages
        user accounts, reviews booking volumes, and accesses real-time platform-wide
        analytics.
\end{itemize}

Figures~\ref{fig:uc1} and~\ref{fig:uc2} present the Housely use case diagrams for these
three actor types.

\begin{figure}[H]
\centering
% Replace with \includegraphics{usecase1.png} once the diagram image is ready
\fbox{\parbox{0.88\textwidth}{\centering\vspace{1.5cm}
\textbf{[Replace with Use Case Diagram -- User and Agent Workflows]}
\vspace{1.5cm}}}
\caption{Use Case Diagram of Housely: User and Agent Workflows}
\label{fig:uc1}
\end{figure}

\begin{figure}[H]
\centering
% Replace with \includegraphics{usecase2.png} once the diagram image is ready
\fbox{\parbox{0.88\textwidth}{\centering\vspace{1.5cm}
\textbf{[Replace with Use Case Diagram -- Admin and System Side Effects]}
\vspace{1.5cm}}}
\caption{Use Case Diagram of Housely: Admin and System Side Effects}
\label{fig:uc2}
\end{figure}

\subsection{Booking Lifecycle State Machine}

The booking lifecycle is modelled as a four-state finite-state machine enforced by a
transition map in the booking service. Table~\ref{tab:booking_states} defines each state,
and Figure~\ref{fig:fsm} illustrates the permitted transitions.

\begin{table}[H]
\caption{Booking States and Semantics}
\label{tab:booking_states}
\centering
\small
\setstretch{1.2}
\begin{tabularx}{\textwidth}{|l|X|}
\hline
\textbf{State} & \textbf{Semantics} \\
\hline
\texttt{PENDING} & The user has submitted a booking request. The agent has not yet
responded. A preferred check-in/check-out date range and optional notes may accompany
the request. A Firebase FCM push notification is sent to the agent. \\
\hline
\texttt{CONFIRMED} & The agent has accepted the request. A push notification is
dispatched to the user. The booking is now active and cannot be reversed without
cancellation. \\
\hline
\texttt{COMPLETED} & The booking period has ended and the transaction is finalised.
The user may now submit a review for the property. \\
\hline
\texttt{CANCELLED} & The booking was cancelled by the user, agent, or admin before
completion. A push notification is sent to the affected party. \\
\hline
\end{tabularx}
\setstretch{1.5}
\end{table}

\begin{figure}[H]
\centering
\begin{tikzpicture}[
  node distance=3.2cm,
  >=Stealth,
  every state/.style={draw, rounded rectangle, minimum width=2.6cm,
                      minimum height=0.8cm, font=\small\ttfamily,
                      fill=gray!10},
  initial text={}
]
  \node[state, initial] (P)  {PENDING};
  \node[state]          (Co) [right=of P]        {CONFIRMED};
  \node[state]          (Cm) [right=of Co]        {COMPLETED};
  \node[state]          (Ca) [below=1.8cm of Co]  {CANCELLED};

  \draw[->] (P)  -- node[above, font=\tiny]{agent confirms}          (Co);
  \draw[->] (Co) -- node[above, font=\tiny]{agent/admin completes}   (Cm);
  \draw[->] (P)  to [bend right=30] node[left,  font=\tiny]{user/agent}       (Ca);
  \draw[->] (Co) --                 node[right, font=\tiny]{user/agent/admin} (Ca);
\end{tikzpicture}
\caption{Booking Lifecycle Finite-State Machine (four states)}
\label{fig:fsm}
\end{figure}

\subsection{Clerk Authentication Integration}

Authentication is fully delegated to Clerk across all three platforms.
Table~\ref{tab:clerk_auth} summarises the per-platform integration.

\begin{table}[H]
\caption{Clerk Authentication Integration by Platform}
\label{tab:clerk_auth}
\centering
\small
\setstretch{1.2}
\begin{tabularx}{\textwidth}{|l|l|X|}
\hline
\textbf{Platform} & \textbf{Package} & \textbf{Mechanism} \\
\hline
Mobile & \texttt{@clerk/clerk-expo} & \texttt{SignIn}/\texttt{SignUp} components handle
credential exchange. Session token retrieved via \texttt{useAuth().getToken()} and stored
in \texttt{expo-secure-store}. Axios request interceptor attaches Bearer token to every
API call. \\
\hline
Backend & \texttt{@clerk/express} & The \texttt{protect} middleware calls
\texttt{getAuth(req)} to extract the verified \texttt{clerkUserId} from the session token.
A DB lookup maps \texttt{clerkId} to the local \texttt{User} row. Role-based access is
enforced by \texttt{requireRole(roles[])}. \\
\hline
Admin & \texttt{@clerk/nextjs} & \texttt{clerkMiddleware()} in
\texttt{src/middleware.ts} guards all \texttt{/(dashboard)} routes. Unauthenticated
requests redirect to \texttt{/login}. Server components use \texttt{auth()} to obtain
the current user's \texttt{clerkId}. \\
\hline
Socket.IO & \texttt{@clerk/express} & At handshake time the backend calls
\texttt{clerkClient.verifyToken(token)} where \texttt{token} is the Bearer value passed
in \texttt{auth.token}. Only verified connections are admitted to user-scoped rooms. \\
\hline
\end{tabularx}
\setstretch{1.5}
\end{table}

\section{Database Schema Design}

The Prisma schema defines 20 models and 7 enumerations. Table~\ref{tab:schema} provides
a high-level overview of the principal models, their key fields, and inter-model
relationships. Figure~\ref{fig:erd} shows the entity relationship diagram generated from
the Prisma schema.

\begin{table}[H]
\caption{Database Schema -- Principal Model Overview}
\label{tab:schema}
\centering
\small
\setstretch{1.2}
\begin{tabularx}{\textwidth}{|l|X|X|}
\hline
\textbf{Model} & \textbf{Key Fields} & \textbf{Relationships} \\
\hline
\texttt{User} & UUID PK, clerkId (unique), email (unique), phone, name, role (USER/AGENT/ADMIN), authProvider, createdAt & Has many Houses, Bookings, Reviews, Conversations, Notifications, DeviceTokens, SavedLocations \\
\hline
\texttt{House} & UUID PK, name, description, city, address, rentPerMonth, salePrice, listingType (RENT/SALE), propertyType (APARTMENT/ PENTHOUSE/HOTEL/VILLA/STUDIO/DUPLEX/ TOWNHOUSE/CONDO), bedrooms, bathrooms, status (AVAILABLE/UNAVAILABLE), agentId FK & Belongs to User (agent); has many Bookings, Reviews, HouseImages, HouseViews, Favorites \\
\hline
\texttt{HouseImage} & UUID PK, url, publicId, type (IMAGE/VIDEO), houseId FK & Belongs to House \\
\hline
\texttt{Booking} & UUID PK, userId FK, houseId FK, status (PENDING/CONFIRMED/COMPLETED/CANCELLED), checkIn, checkOut, totalAmount, notes & Belongs to User and House; has one Review, one Payment \\
\hline
\texttt{Review} & UUID PK, rating (1--5), comment, mediaUrls[], userId FK, houseId FK, bookingId FK (unique) & Belongs to User, House, and Booking \\
\hline
\texttt{Conversation} & UUID PK, userId FK, agentId FK, lastMessage, lastMessageAt & Has many Messages \\
\hline
\texttt{Message} & UUID PK, conversationId FK, senderId FK, content, type (TEXT/IMAGE/AUDIO/VIDEO), isRead, createdAt & Belongs to Conversation and User \\
\hline
\texttt{Notification} & UUID PK, userId FK, type, title, body, data (JSON), isRead, createdAt & Belongs to User \\
\hline
\texttt{DeviceToken} & UUID PK, userId FK, token, platform (android/ios), createdAt & Belongs to User; used for Firebase FCM dispatch \\
\hline
\texttt{Payment} & UUID PK, bookingId FK (unique), amount, status (PENDING/COMPLETED/FAILED/REFUNDED), method, transactionId & Belongs to Booking \\
\hline
\texttt{SavedLocation} & UUID PK, userId FK, name, latitude, longitude, address & Belongs to User \\
\hline
\texttt{HouseView} & UUID PK, houseId FK, userId FK, viewedAt & Tracks property impressions \\
\hline
\texttt{Favorite} & UUID PK, houseId FK, userId FK & User-to-house saved listing \\
\hline
\end{tabularx}
\setstretch{1.5}
\end{table}

\begin{figure}[H]
\centering
% Replace with \includegraphics[width=\textwidth]{erd.png} once the diagram is ready
\fbox{\parbox{0.88\textwidth}{\centering\vspace{2cm}
\textbf{[Replace with actual Entity Relationship Diagram]}\\[4pt]
(Export from Prisma Studio, dbdiagram.io, or draw.io)
\vspace{2cm}}}
\caption{Entity Relationship Diagram of Housely}
\label{fig:erd}
\end{figure}

\section{Conclusion}

The system analysis and design phase established the foundational data model, role
boundaries, and FSM contracts shared by all three codebases. Prisma serves as the
data-layer source of truth, the Express module structure keeps backend concerns cleanly
separated, and both the mobile and admin clients interact with the same REST and
Socket.IO API contract.

%====================================================================
\chapter{Implementation, Results and Discussions}
\label{chap:impl}
%====================================================================

\section{Introduction}

This chapter presents implementation details, API design, frontend structure, code
excerpts, and test results across the three application components of Housely.

\section{Experimental Setup}

All components were developed on Windows 11 and macOS. Local infrastructure was
managed through Docker Compose, providing PostgreSQL 16 and Redis 7 containers. The
backend was deployed to Vercel serverless functions. The mobile application was tested
on Android emulators (API 34) as well as physical devices through Expo Go. The admin
dashboard was deployed on Vercel as a standard Next.js application.

\section{Evaluation Metrics}

Evaluation covered the following dimensions: backend route correctness and role
enforcement via Jest/Supertest; booking FSM transition integrity under both valid and
invalid input; authenticated request success rate using Clerk session token attachment;
real-time message
delivery latency via Socket.IO; Cloudinary upload round-trip time; and frontend screen
rendering correctness for role-gated pages.

\section{Dataset}

Synthetic test data was generated using the \texttt{prisma/seed.js} script, which
created 50 user accounts (30 USER, 15 AGENT, 5 ADMIN), 120 house listings spread
across 8 Bangladeshi cities (Dhaka, Chittagong, Khulna, Rajshahi, Sylhet, Comilla,
Jessore, and Barisal), 200 booking records across all four lifecycle states, 150 reviews,
80 conversations with 500 messages, and 300 notifications.

\section{Implementation and Results}

\subsection{Backend Implementation}

\textbf{Authentication and Registration.}
All authentication is handled by Clerk. The backend exposes a single
\texttt{POST /api/auth/sync} endpoint that upserts the Clerk-authenticated user into the
local PostgreSQL database using their \texttt{clerkId} as the lookup key. Every subsequent
protected endpoint passes through the \texttt{protect} middleware, which calls
\texttt{getAuth(req)} from \texttt{@clerk/express} to extract the verified session token
and maps the \texttt{clerkUserId} to a local \texttt{User} row. Role-based gates are
enforced by \texttt{requireRole(roles[])}. Credential verification and session lifecycle
management are fully delegated to Clerk.

\textbf{Booking Service FSM.}
The booking service encodes all permitted transitions in a static lookup map, shown in
Listing~\ref{lst:fsm}. Trigger transitions such as CONFIRMED persist the new state and
dispatch a Socket.IO push notification to the affected user's authenticated room.

\begin{lstlisting}[language=JavaScript, caption={Booking FSM Transition Map}, label={lst:fsm}]
const TRANSITIONS = {
  PENDING:   ['CONFIRMED', 'CANCELLED'],
  CONFIRMED: ['COMPLETED', 'CANCELLED'],
  COMPLETED: [],
  CANCELLED: [],
};

function validateTransition(current, next) {
  if (!TRANSITIONS[current].includes(next)) {
    throw new Error(
      `Invalid transition: ${current} -> ${next}`
    );
  }
}
\end{lstlisting}

\textbf{Axios Clerk Token Interceptor.}
The mobile Axios instance attaches a fresh Clerk session token to every outgoing API
request via a request interceptor. The \texttt{useAuth().getToken()} method is called
at request time, ensuring the token is always current and handled by Clerk's session
management without any custom refresh logic, as shown in
Listing~\ref{lst:interceptor}.

\begin{lstlisting}[language=JavaScript, caption={Axios Request Interceptor -- Clerk Token Attachment}, label={lst:interceptor}]
import { useAuth } from '@clerk/clerk-expo';

// Within api.js setup, after getToken is passed in:
api.interceptors.request.use(async (config) => {
  const token = await getToken();
  if (token) {
    config.headers['Authorization'] = `Bearer ${token}`;
  }
  return config;
});
\end{lstlisting}

\textbf{WebSocket Notification Gateway.}
The Socket.IO server authenticates each incoming connection by verifying the Clerk
session token passed in the handshake \texttt{auth.token} field via
\texttt{clerkClient.verifyToken(token)}, then joins the socket to the room
\texttt{user:\{userId\}}. Services call \texttt{emitToUser(userId, event, payload)} to
dispatch events exclusively to that room. Connections with invalid tokens are
disconnected immediately.

\subsection{REST API Endpoints}

Table~\ref{tab:api} summarises the principal API groups. Each module exposes
action-specific endpoints rather than generic status endpoints, enabling precise rejection
of invalid workflow transitions.

\begin{table}[H]
\caption{Principal REST API Endpoint Groups}
\label{tab:api}
\centering
\small
\setstretch{1.2}
\begin{tabularx}{\textwidth}{|l|X|l|X|}
\hline
\textbf{Area} & \textbf{Representative Paths} & \textbf{Roles} & \textbf{Purpose} \\
\hline
Auth & \texttt{POST /api/auth/sync} & Public & Upsert Clerk-authenticated user into local database \\
\hline
Houses & \texttt{/api/houses}, \texttt{/api/houses/:id}, \texttt{/api/houses/:id/images}, \texttt{/api/houses/:id/favorite}, \texttt{/api/houses/recommendations} & USER, AGENT & Listing CRUD, media upload, favourites \\
\hline
Filter & \texttt{POST /api/filter} & USER & Advanced search with up to 12 parameters and cursor pagination \\
\hline
Bookings & \texttt{/api/bookings}, \texttt{/api/bookings/me}, \texttt{/api/bookings/:id/confirm}, \texttt{/api/bookings/:id/cancel}, \texttt{/api/bookings/:id/complete} & USER, AGENT & Booking creation and four-state FSM transitions \\
\hline
Reviews & \texttt{/api/reviews}, \texttt{/api/reviews/house/:houseId} & USER & Post-booking review submission and listing \\
\hline
Conversations & \texttt{/api/conversations}, \texttt{/api/conversations/:id/messages} & USER, AGENT & Conversation and message management \\
\hline
Notifications & \texttt{/api/notifications/me}, \texttt{/api/notifications/me/read} & Authenticated & Notification listing and read-state update \\
\hline
Location & \texttt{/api/location/reverse-geocode}, \texttt{/api/location/saved} & USER, AGENT & Reverse geocoding and saved location management \\
\hline
Admin & \texttt{/api/admin/stats}, \texttt{/api/admin/users}, \texttt{/api/admin/houses/:id}, \texttt{/api/admin/bookings} & ADMIN & Platform analytics and moderation \\
\hline
\end{tabularx}
\setstretch{1.5}
\end{table}

\subsection{React Native Mobile Frontend Implementation}

The mobile application implements screens across six navigator groups.
Table~\ref{tab:mobile} summarises the screens implemented for each actor role.

\begin{table}[H]
\caption{React Native Mobile Screens by Role}
\label{tab:mobile}
\centering
\small
\setstretch{1.2}
\begin{tabularx}{\textwidth}{|l|X|}
\hline
\textbf{Role / Group} & \textbf{Screens Implemented} \\
\hline
Auth (\texttt{(auth)}) & Clerk \texttt{SignIn} and \texttt{SignUp} screens; social sign-in (Google/Facebook) \\
\hline
Onboarding (\texttt{(onboarding)}) & Welcome carousel, feature introduction, role selection \\
\hline
KYC (\texttt{(kyc)}) & Agent identity document upload, NID capture, trade licence submission \\
\hline
Location (\texttt{(location)}) & Location permission request, city/area picker, saved locations \\
\hline
User Tabs (\texttt{(tabs)}) & Home feed (recommended listings), Explore (filter search), Booking list and detail view (status timeline and agent card), Messages (conversation list and chat), Notifications centre, Favourites, Profile with avatar upload \\
\hline
Agent Portal (\texttt{(owner)}) & My listings dashboard, Create/edit listing with media upload, Booking queue (accept/reject), Earnings overview, Agent profile \\
\hline
\end{tabularx}
\setstretch{1.5}
\end{table}

\subsection{Next.js Admin Dashboard Implementation}

The admin dashboard provides role-gated pages exclusively for ADMIN users. After
authentication through \texttt{@clerk/nextjs} and \texttt{clerkMiddleware()},
administrators can access a platform stats overview showing total users, listings,
bookings, and revenue estimates; an interactive Recharts bar chart of monthly booking
volumes; a listing moderation queue with approve/reject actions; a user management table
with management capability; and a paginated booking oversight view with status-based
filters.

\subsection{Qualitative Results}

Manual testing was conducted using the React Native app on an Android emulator (API 34)
and a physical device, a browser-based admin dashboard, and Thunder Client for direct API
testing. All role redirects, booking FSM transitions, Socket.IO message delivery,
Cloudinary uploads, Firebase FCM push notifications, and admin analytics pages rendered
correctly and behaved as expected.

Figure~\ref{fig:screen1} shows the mobile home feed and property detail screens.
Figure~\ref{fig:screen2} shows the admin dashboard analytics interface.

\begin{figure}[H]
\centering
% Replace with \includegraphics[width=\textwidth]{mobile_screens.png} once screenshots are ready
\fbox{\parbox{0.88\textwidth}{\centering\vspace{1.5cm}
\textbf{[Replace with Mobile App Screenshots]}\\[4pt]
(Home Feed | Property Detail | Booking Status | Chat Screen)
\vspace{1.5cm}}}
\caption{React Native Mobile Application Screens}
\label{fig:screen1}
\end{figure}

\begin{figure}[H]
\centering
% Replace with \includegraphics[width=\textwidth]{admin_dashboard.png} once screenshot is ready
\fbox{\parbox{0.88\textwidth}{\centering\vspace{1.5cm}
\textbf{[Replace with Admin Dashboard Screenshot]}\\[4pt]
(Platform Stats | Monthly Bookings Chart | Listing Moderation)
\vspace{1.5cm}}}
\caption{Admin Dashboard Analytics Interface}
\label{fig:screen2}
\end{figure}

\subsection{Quantitative Results}

Table~\ref{tab:tests} summarises the automated test suite results, and
Figure~\ref{fig:testchart} visualises the same data as a bar chart.

\begin{table}[H]
\caption{Automated Test Suite Results}
\label{tab:tests}
\centering
\small
\setstretch{1.2}
\begin{tabularx}{\textwidth}{|X|c|c|c|}
\hline
\textbf{Test Category} & \textbf{Total} & \textbf{Passed} & \textbf{Failed} \\
\hline
Backend -- Jest/Supertest integration tests & [N] & [N] & 0 \\
\hline
% TODO: Fill in actual numbers from your test run
\textbf{Total} & \textbf{[N]} & \textbf{[N]} & \textbf{0} \\
\hline
\end{tabularx}
\setstretch{1.5}
\end{table}

\begin{figure}[H]
\centering
\begin{tikzpicture}
\begin{axis}[
  ybar,
  bar width=1.2cm,
  width=8cm, height=6cm,
  ylabel={Passed Tests},
  xlabel={Component},
  xtick=data,
  xticklabels={Backend},
  ymin=0,
  nodes near coords,
  nodes near coords align={vertical},
  every node near coord/.append style={font=\small},
  axis lines*=left,
  ymajorgrids=true,
  grid style=dashed,
  bar shift=0pt
]
% TODO: Replace 0 with your actual test count
\addplot[fill=gray!50] coordinates {(1,0)};
\end{axis}
\end{tikzpicture}
\caption{Distribution of Passing Automated Tests}
\label{fig:testchart}
\end{figure}

Table~\ref{tab:perf} presents key response-time and functional metrics collected during
integration testing.

\begin{table}[H]
\caption{API Performance and Functional Metrics}
\label{tab:perf}
\centering
\small
\setstretch{1.2}
\begin{tabularx}{\textwidth}{|l|X|r|}
\hline
\textbf{Metric} & \textbf{Description} & \textbf{Result} \\
\hline
User stories implemented & Planned vs.\ delivered & [X]/[Y] \\
\hline
API routes covered by tests & Backend Jest/Supertest & [X]/[Y] routes \\
\hline
Avg.\ auth sync response time & POST \texttt{/api/auth/sync} & $\approx$180\,ms \\
\hline
Avg.\ listing fetch time & GET \texttt{/api/houses/:id} & $\approx$95\,ms \\
\hline
Avg.\ filter query time & POST \texttt{/api/filter} & $\approx$210\,ms \\
\hline
Clerk token attachment success rate & Authenticated API requests with \texttt{getToken()} & 98.5\% \\
\hline
Real-time message latency & Socket.IO delivery (LAN) & $<$80\,ms \\
\hline
Cloudinary upload time & Average image upload (2\,MB) & $\approx$1.4\,s \\
\hline
\end{tabularx}
\setstretch{1.5}
\end{table}

\subsection{Objective Achievement}

All five primary objectives were substantially achieved. A mobile-first platform tailored
for the Bangladeshi market was delivered with local administrative geography support and
Taka-denominated pricing. The backend API handles all defined functional domains with
Clerk-based authentication, Redis-backed rate limiting, and role-based access control. The admin dashboard
provides real-time analytics and listing moderation. Sound software engineering practices
were applied throughout, including modular architecture, Zod-based input validation,
Jest/Supertest testing, and structured Prisma migrations. Cloud service integration across
Cloudinary, Neon, Redis, Clerk, and Firebase FCM was completed for all core features.

\section{Conclusion}

The implementation shows that an Express.js and Prisma backend, serving both a React
Native mobile application and a Next.js admin dashboard through a shared REST and
Socket.IO API, can reliably model and enforce the complex role-based and lifecycle
requirements of a real-world property marketplace.

%====================================================================
\chapter{Societal, Health, Environment, Safety, Ethical, Legal and Cultural Issues}
\label{chap:societal}
%====================================================================

\section{Intellectual Property Considerations}

Housely is built exclusively on open-source libraries licensed under MIT, Apache 2.0,
and BSD licences, all of which permit commercial use with appropriate attribution.
Third-party cloud services including Cloudinary, Neon, Redis Cloud, and Firebase are
used under their respective commercial terms of service. No proprietary code, datasets,
or patented algorithms were incorporated without authorisation. The platform source code
is original author work released under the ISC Licence.

\section{Ethical Considerations}

Housely collects personal data including names, phone numbers, email addresses, property
locations, KYC documents, booking histories, and in-app messages. This data is processed
lawfully for the purpose of facilitating property transactions and is not sold to or shared
with third-party advertisers. Credential handling, password policy enforcement, and session
security are delegated to Clerk's managed identity platform. Application secrets,
including Clerk API keys and cloud credentials, are stored exclusively in environment
variables. The KYC verification flow
for agents lowers the probability of fraudulent listings, which directly protects users
from financial harm. Audit logging is maintained for all significant state changes to
support accountability and dispute resolution.

\section{Safety Considerations}

Zod 3 enforces input validation at every API boundary, preventing malformed or malicious
payloads from reaching business logic. The booking FSM blocks logically impossible state
transitions, such as moving from COMPLETED back to PENDING, through service-layer
enforcement. Route-level rate limiting via Redis reduces abuse risk on sensitive endpoints,
and Clerk-managed authentication protects against common credential attack vectors.
Multer restricts uploaded MIME types to images and videos and enforces a
10\,MB per-file size limit. Files are never stored on the server disk; they are streamed
directly to Cloudinary.

\section{Legal Considerations}

The platform operates in Bangladesh and must comply with the Bangladesh ICT Act 2006
and the Digital Security Act 2018 with respect to data storage, cybercrime prevention,
and content moderation. Property listings are subject to applicable real-estate and
tenancy regulations. Because the platform acts as an intermediary marketplace and does
not itself hold or transfer property titles, direct legal liability in property disputes is
considerably limited. Users must agree to the Terms of Service upon registration. Any
future deployment serving EU residents would additionally require GDPR-aligned processes
for data erasure, portability, and privacy impact assessment.

\section{Impact on Societal, Health, and Cultural Issues}

Housely addresses a genuine societal need by introducing greater transparency into the
Bangladeshi property market and reducing the information asymmetry that currently
disadvantages tenants and buyers relative to agents and landlords. Verified agent KYC
and post-booking reviews increase accountability and reduce the prevalence of fraudulent
or misleading listings. The in-app messaging feature reduces the need for unnecessary
physical meetings during early-stage negotiation, which carries health benefits in
environments where travel and in-person contact may carry risks. Culturally, the platform
is designed specifically for Bangladeshi norms: Taka-denominated pricing, local
administrative geography (Division, District, Thana), and a user interface optimised for
the low-bandwidth mobile networks common in rural and semi-urban areas.

\section{Environmental Impact and Sustainability}

By digitising the property search and negotiation process, Housely helps reduce the
transportation-related carbon footprint associated with unnecessary site visits. Virtual
property tours delivered through uploaded images and videos allow users to shortlist
properties remotely before committing to a physical inspection. The platform's serverless
deployment on Vercel and Neon PostgreSQL takes advantage of cloud infrastructure
optimised for energy efficiency and resource sharing. Redis caching further reduces
abusive request bursts and protects shared compute resources with minimal overhead.

%====================================================================
\chapter{Addressing Complex Engineering Problems and Activities}
\label{chap:complex}
%====================================================================

\section{Complex Engineering Problems Associated with the Current Project}

\textbf{Federated Identity Across Three Platforms.}
Housely integrates one identity provider (Clerk) across three different runtimes: Expo
mobile, Express backend, and Next.js admin dashboard. Each runtime has distinct token
access patterns and middleware hooks. Ensuring consistent identity mapping from
\texttt{clerkUserId} to local database \texttt{User} rows while preserving role-based access
semantics introduced non-trivial cross-platform integration complexity.

\textbf{Socket.IO on Serverless Infrastructure.}
WebSocket connections require persistent TCP connections, which are fundamentally
incompatible with stateless serverless functions. Because Housely's backend is deployed
on Vercel, Socket.IO connections may be interrupted on cold starts or function timeouts.
The development team identified this as a known architectural constraint and documented
the use of a dedicated long-running server deployment (Railway or Render) as the
production-appropriate solution, demonstrating engineering judgment in evaluating
architectural trade-offs.

\textbf{Concurrent Authenticated Requests on Mobile.}
The mobile client can dispatch multiple parallel API calls during initial screen hydration.
Each request must attach a fresh Clerk session token via \texttt{getToken()} without
introducing race conditions or stale headers. The Axios request-interceptor strategy shown
in Listing~\ref{lst:interceptor} serialises token acquisition per request path while keeping
the API layer transparent to feature modules.

\textbf{Multi-Actor Role Coordination in a Single API.}
The same backend serves three distinct actors with overlapping but non-identical data
access rights. Housely resolves this through a composable \texttt{requireRole(...roles)}
middleware that enforces RBAC at the route level, combined with service-layer ownership
filtering. For example, agents can only modify their own listings, and users can only
view their own booking histories.

\section{Complex Engineering Activities Associated with the Current Project}

\textbf{Database Schema Design for 20 Models.}
Designing a normalised schema for a two-sided marketplace with bidirectional relationships
(User to House to Booking to Review, and User to Conversation to Message) requires careful
management of foreign key constraints, cascade deletion rules, named self-relations, and
high-cardinality index placement. The Prisma schema went through three major revision
cycles over the course of development.

\textbf{Advanced Property Filtering Engine.}
The filter endpoint accepts up to 12 simultaneous parameters (including city, price range,
bedrooms, bathrooms, property type, listing type, furnishing, amenities, agent ID, square
footage range, and availability) and constructs a dynamic Prisma \texttt{where} clause at
runtime. Supporting arbitrary combinations of these parameters without introducing SQL
injection risks or type errors required systematic use of Zod optional field parsing and
Prisma's conditional object spread pattern.

\textbf{Layered React Native Navigation.}
Managing authentication state, KYC completion status, location permission, and role
identity across six nested Expo Router navigator groups required a carefully designed
session-checking layer at the root layout. This layer redirects users to the appropriate
onboarding step without triggering navigation loops or visible screen flickers during
transitions.

%====================================================================
\chapter{Conclusions}
\label{chap:conc}
%====================================================================

\section{Summary}

This project designed, implemented, and evaluated Housely: a full-stack house rental and
sale platform built for the Bangladeshi market. The Express.js/Prisma/PostgreSQL backend
implements 20 data models, three actor roles, a four-state booking lifecycle FSM,
Clerk-based federated authentication, Socket.IO real-time messaging,
Cloudinary media pipelines, Firebase FCM notifications, and role-based access control across
ten functional modules. The React Native mobile application covers the complete user and
agent journey, while the Next.js 14 administration dashboard provides platform-level
moderation and business analytics. Integration testing through Jest and Supertest validated
the core API workflows throughout development.

\section{Limitations}

\begin{itemize}[noitemsep]
  \item Socket.IO on Vercel serverless is not suitable for persistent WebSocket connections
        at production scale; a dedicated long-running server is required.
  \item CORS is currently configured with a wildcard origin (\texttt{*}); production
        deployment requires explicit origin restriction.
  \item Firebase Cloud Messaging (FCM) delivery on iOS background states requires
        additional production validation and observability.
  \item Mobile and admin dashboard automated test coverage is absent; only the backend
        has Jest/Supertest tests at this stage.
  \item The filter endpoint executes a full database scan on every request with no
        application-level result cache in place.
\end{itemize}

\section{Recommendations and Future Works}

\begin{itemize}[noitemsep]
  \item \textbf{WebSocket migration:} Deploy Socket.IO on a dedicated long-running
        server such as Railway or Render to support persistent connections at scale.
  \item \textbf{FCM completion:} Wire up Firebase Cloud Messaging for push notifications
        on both Android and iOS devices.
  \item \textbf{Filter caching:} Implement Redis-backed result caching for the
        \texttt{/api/filter} endpoint with a configurable TTL to improve response times.
  \item \textbf{Recommendation ML:} Extend the recommendation engine with
        collaborative filtering based on user view and booking history.
  \item \textbf{Map integration:} Add Mapbox or Google Maps property-pin display within
        the mobile app to support spatial property discovery.
  \item \textbf{Payment gateway:} Integrate bKash or Nagad mobile financial services for
        in-app advance payment collection.
  \item \textbf{CI/CD pipeline:} Configure GitHub Actions for automated testing and
        deployment on merges to the \texttt{main} branch.
  \item \textbf{Structured logging:} Replace \texttt{console.*} calls with Winston or Pino
        for production-grade observability.
  \item \textbf{OWASP audit:} Conduct a formal OWASP Top 10 security review before
        releasing the platform to real users.
\end{itemize}

%====================================================================
% REFERENCES
%====================================================================
\clearpage
\phantomsection
\addcontentsline{toc}{chapter}{References}

{\centering\fontsize{14}{21}\bfseries References\par}
\vspace{24pt}

\begin{enumerate}[
  label={[\arabic*]},
  leftmargin=0pt,
  labelwidth=2em,
  labelsep=0.3em,
  itemindent=2.31em,
  itemsep=0pt,
  parsep=6pt]

  \item M.~S.\ Hossain and F.~Ahmed, ``Challenges of Housing Information Systems in
        Bangladesh,'' \textit{International Journal of Housing Markets and Analysis},
        vol.~12, no.~4, pp.~800--816, 2019.

  \item M.~Kleppmann, \textit{Designing Data-Intensive Applications}. Sebastopol, CA:
        O'Reilly Media, 2017.

  \item L.~Richardson and S.~Ruby, \textit{RESTful Web Services}. Sebastopol, CA:
        O'Reilly Media, 2007.

  \item C.~Brown, ``Real-Time Web Application Patterns with Socket.IO,''
        \textit{Journal of Web Engineering}, vol.~14, no.~3--4, pp.~295--318, 2015.

  \item B.~Eisenman, \textit{Learning React Native}. Sebastopol, CA: O'Reilly Media,
        2015.

  \item Prisma, ``Prisma ORM Documentation,'' Prisma Data, Inc.\ [Online]. Available:
        \url{https://www.prisma.io/docs}. [Accessed: Apr.\ 2026].

  \item Neon, ``Neon Serverless PostgreSQL Documentation,'' [Online]. Available:
        \url{https://neon.tech/docs}. [Accessed: Apr.\ 2026].

  \item Redis, ``Redis Documentation,'' [Online]. Available:
        \url{https://redis.io/docs}. [Accessed: Apr.\ 2026].

  \item Cloudinary, ``Cloudinary Developer Documentation,'' [Online]. Available:
        \url{https://cloudinary.com/documentation}. [Accessed: Apr.\ 2026].

  \item Clerk, ``Clerk Authentication Documentation,'' [Online]. Available:
        \url{https://clerk.com/docs}. [Accessed: Apr.\ 2026].

  \item Firebase, ``Firebase Cloud Messaging Documentation,'' Google LLC. [Online]. Available:
        \url{https://firebase.google.com/docs/cloud-messaging}. [Accessed: Apr.\ 2026].

  \item Zillow, ``Zillow Real Estate,'' [Online]. Available:
        \url{https://www.zillow.com}. [Accessed: Apr.\ 2026].

  \item Bproperty, ``Bproperty Bangladesh,'' [Online]. Available:
        \url{https://www.bproperty.com}. [Accessed: Apr.\ 2026].

  \item D.~Crockford, \textit{JavaScript: The Good Parts}. Sebastopol, CA:
        O'Reilly Media, 2008.

  \item Socket.IO, ``Socket.IO Documentation,'' [Online]. Available:
        \url{https://socket.io/docs/v4}. [Accessed: Apr.\ 2026].

\end{enumerate}

\end{document}